\section{Work the four-call state trace}\label{utf8-x-trace}
At each \texttt{chunk\_ready} marker, main has published the call result and the
scanner's state is coherent. The table uses decimal counts and offsets, with the
next-byte range written in hexadecimal. Need zero means that no continuation is
pending; the stored ordinary bounds then provide the canonical idle values.

\Needspace{12\baselineskip}
\begin{longtable}{rrrrrlr}
  \toprule Call & New bytes & Total & Count & Need & Next hex        & Start  \\\midrule
  0             & 2         & 2     & 1     & 2    & \texttt{80--bf} & 1      \\
  1             & 3         & 5     & 2     & 3    & \texttt{90--bf} & 4      \\
  2             & 2         & 7     & 2     & 1    & \texttt{80--bf} & 4      \\
  3             & 4         & 11    & 5     & 0    & \texttt{80--bf} & \(-1\) \\\bottomrule
\end{longtable}
The first three statuses are zero. The fourth is one because its final flag is
set and no sequence remains pending. Every error word is minus one. Count does not
change in call two, although that call accepts two bytes. The lead start four is
retained until the last byte of that encoding arrives in call three.

\begin{MARSWorkedExample}{A byte error and an end error}
  Change input byte five from \texttt{9f} to \texttt{8f}. The third call begins with a
  byte below the retained \texttt{f0} minimum \texttt{90}. It returns minus one and
  zero accepted bytes. Stored total is five, count two, need three, start four, and
  error five. The fourth call returns minus three without consuming anything, while
  the stored minus-one status and first diagnostic remain unchanged.

  Restore the input and instead make the fourth chunk empty while retaining its
  final flag. That call returns minus two and zero accepted bytes. Total remains
  seven, count two, need one, start four, and error seven. The mapped \texttt{82} at
  the supplied input pointer is outside this zero-byte chunk and must not be read.
\end{MARSWorkedExample}
